\documentclass[sigconf,balance=false]{webmedia}

\AtBeginDocument{%
  \providecommand\BibTeX{{%
    \normalfont B\kern-0.5em{\scshape i\kern-0.25em b}\kern-0.8em\TeX}}}
\settopmatter{printacmref=false, printfolios=false}
\setevent{webmedia}
\proceedingsDetails[PUC MG]{Otimização de Sistemas}{2026}{Belo Horizonte/MG, Brazil} 
% \ISSN{2966-2753}

%\setevent{ctd}
%\proceedingsDetails[CTD’2026]{VIII Concurso de Teses e Disserta{\c c}{\~o}es. Anais Estendidos do XXXII Simpósio Brasileiro de Sistemas Multimídia e Web}{2026}{Lavras/MG, Brasil} 
%\ISSN{2596-1683}

%\setevent{ctic}
%\proceedingsDetails[CTIC’2026]{VI Concurso de Trabalhos de Inicia{\c c}{\~a}o Cient{\'i}fica. Anais Estendidos do XXXII Simpósio Brasileiro de Sistemas Multimídia e Web}{2026}{Lavras/MG, Brasil} 
%\ISSN{2596-1683}

%\setevent{wfa}
%\proceedingsDetails[WFA’2026]{XXV Workshop de Ferramentas e Aplica{\c c}{\~o}es. Anais Estendidos do XXXII Simpósio Brasileiro de Sistemas Multimídia e Web}{2026}{Lavras/MG, Brasil} 
%\ISSN{2596-1683}

%\setevent{w4e}
%\proceedingsDetails[W4E’2026]{V WebMedia for Everyone. Anais Estendidos do XXXII Simpósio Brasileiro de Sistemas Multimídia e Web}{2026}{Lavras/MG, Brasil} 
%\ISSN{2596-1683}

%\setevent{wrsl}
%\proceedingsDetails[WRSL+’2026]{II Workshop de Revisões Sistemáticas de Literatura em Sistemas Multimídia e Web. Anais Estendidos do XXXII Simpósio Brasileiro de Sistemas Multimídia e Web}{2026}{Lavras/MG, Brasil} 
%\ISSN{2596-1683}

%\setevent{BrWeb3}
%\proceedingsDetails[BrWeb3’2026]{II Workshop Brasileiro de Sistemas Web3. Anais Estendidos do XXXII Simpósio Brasileiro de Sistemas Multimídia e Web}{2026}{Lavras/MG, Brasil} 
%\ISSN{2596-1683}

%\setevent{wtvdi}
%\proceedingsDetails[WTVDI’2026]{IX Workshop Futuro da TV Digital Interativa. Anais Estendidos do XXXII Simpósio Brasileiro de Sistemas Multimídia e Web}{2026}{Lavras/MG, Brasil}
%\ISSN{2596-1683}

%% end of the preamble, the start of the body of the document source.

% extra
\usepackage[utf8]{inputenc}
\usepackage[T1]{fontenc}
% \usepackage[brazil]{babel}
\usepackage[brazilian]{babel}
% newtxmath (carregado por webmedia.cls) já define \Bbbk; amssymb
% tenta defini-lo de novo e aborta a compilação.
\let\Bbbk\relax
\usepackage{amsmath,amssymb}
\usepackage{booktabs}
\usepackage{array}
\usepackage{longtable}
\usepackage{hyperref}
\usepackage{enumitem}
\usepackage{microtype}
\usepackage{graphicx}
\usepackage{float}
\usepackage{balance}
\begin{document}

\title{Feature Engineering e Feature Store}
\subtitle{Pipeline de Dados: Integração, Qualidade e Automação}
\author{Lucas Zegrine Duarte}
\email{1327493@pucminas.edu.br}
\email{lucas.zegrine@gmail.com}
\affiliation{%
  \position{Aluno - 1327493}
  \institution{PUC Minas}
  \department{Pós Graduação Engenharia de Inteligência Artificial e MLOps}
  \city{Belo Horizonte}
  \state{Minas Gerais}
  \country{Brasil}
}

%%
%% The abstract is a short summary of the work to be presented in the
%% article.
% \begin{abstract}
% Este trabalho resolve o Problema 16 de Programação Não Linear, cujo objetivo é distribuir uma demanda total entre quatro máquinas síncronas de modo a minimizar as perdas totais do sistema. O problema é formulado como uma minimização quadrática convexa com uma restrição linear de igualdade e restrições laterais de capacidade. A solução analítica é obtida pelo método dos Multiplicadores de Lagrange, complementado pela verificação das condições de Karush-Kuhn-Tucker para tratar a restrição ativa de limite superior. Em seguida, aplica-se o método de Descida Coordenada, formulado como transferências viáveis de carga entre pares de máquinas, preservando a restrição de balanço em todas as iterações. Os dois procedimentos conduzem a soluções numericamente equivalentes, confirmando a consistência do resultado obtido.
% \end{abstract}


%%
%% Keywords. The author(s) should pick words that accurately describe
%% the work being presented. Separate the keywords with commas.
\keywords{feature engineering, normalização, padronização, seleção de atributos}

%%
%% This command processes the author and affiliation and title
%% information and builds the first part of the formatted document.
\maketitle

\section{Introdução}

\textit{Feature Engineering} é o processo de criar, transformar, selecionar e representar variáveis a partir dos dados brutos, de modo que um modelo de aprendizado de máquina receba entradas numéricas, comparáveis e relevantes para o problema. Os dados originais reúnem categorias, escalas distintas e relações que não aparecem de forma direta nas colunas iniciais. As técnicas a seguir organizam essas transformações.

\begin{itemize}
    \item \textbf{One-Hot Encoding.} Representa uma variável categórica sem ordem natural por um conjunto de colunas binárias, uma para cada categoria. Cada observação recebe o valor 1 na coluna da categoria correspondente e 0 nas demais. Cada coluna identifica uma categoria, o que torna a representação interpretável e compatível com modelos que exigem entradas numéricas, e as categorias permanecem independentes. Com muitas categorias, surge uma coluna por valor distinto: a matriz fica esparsa, o uso de memória cresce e o treinamento de alguns modelos fica mais custoso. Categorias que aparecem só em produção, ausentes no treinamento, pedem tratamento explícito, como o oferecido pelo \texttt{OneHotEncoder} do scikit-learn.

    \item \textbf{Ordinal Encoding.} Representa uma variável categórica ordenada por números inteiros que preservam essa hierarquia. A posição numérica de cada código corresponde à posição relativa da categoria.

    \item \textbf{Normalização.} Reescala cada variável para um intervalo fixo, em geral de 0 a 1, pela transformação Min-Max. Variáveis medidas em unidades ou amplitudes diferentes passam a ocupar a mesma faixa de valores.

    \item \textbf{Padronização.} Recentra cada variável na média zero e a reescala para desvio padrão unitário, pela transformação Z-score. Cada observação passa a ser descrita pelo número de desvios em relação à média da própria variável.

    \item \textbf{Feature Selection.} Escolhe o subconjunto de variáveis mais associadas ao objetivo do modelo. A seleção reduz ruído, redundância e a quantidade de entradas apresentadas ao algoritmo. Os procedimentos se organizam em três categorias:
    \begin{itemize}
        \item \textbf{Filter Methods.} Avaliam cada variável por uma medida estatística calculada sobre os dados, antes de qualquer treinamento. Correlação, informação mútua, teste qui-quadrado e ANOVA são exemplos. O custo permanece baixo e o critério independe do algoritmo usado depois. Cada variável é julgada isoladamente, então interações entre elas ficam fora da seleção.

        \item \textbf{Wrapper Methods.} Percorrem subconjuntos de variáveis e medem cada subconjunto pelo desempenho de um modelo treinado nele. Seleção \textit{forward}, eliminação \textit{backward} e eliminação recursiva de features (RFE) são exemplos. O subconjunto fica ligado ao modelo da busca e reflete combinações entre variáveis. O custo cresce com o número de subconjuntos avaliados.

        \item \textbf{Embedded Methods.} Incorporam a seleção ao próprio treinamento. A regularização L1, como no Lasso, reduz a zero os coeficientes de variáveis pouco úteis, e modelos baseados em árvores atribuem importância às variáveis ao construir as divisões. O critério de seleção e o modelo ficam no mesmo ajuste, com custo próximo ao de um único treinamento.
    \end{itemize}

    \item \textbf{Polynomial Features.} Gera novas variáveis a partir de potências e produtos das variáveis originais. Essas combinações permitem que um modelo linear represente relações curvas e interações entre as entradas.
\end{itemize}

Este relatório aplica essas técnicas às questões da atividade. Cada questão reúne o enunciado, o procedimento adotado e o resultado obtido.

\section{Questão 1}

\subsection{Enunciado}

Um e-commerce possui várias categorias de produtos descritas no conjunto de dados a seguir:
\begin{verbatim}
{
    'Produto': [
        'Eletrônicos',
        'Roupas',
        'Eletrodomésticos',
        'Eletrônicos',
        'Livros',
        'Roupas',
        'Livros'
    ]
}
\end{verbatim}
Aplique a codificação one-hot para representar as categorias de produtos em formato numérico.

\subsection{Solução}

A coluna \texttt{Produto} tem quatro categorias sem ordem: Eletrodomésticos, Eletrônicos, Livros e Roupas. O script monta o \texttt{DataFrame} das sete observações e aplica \texttt{pandas.get\_dummies} com \texttt{dtype=int}. Cada categoria vira uma coluna binária (\texttt{Produto\_<categoria>}): 1 na categoria da linha e 0 nas demais. Cada linha fica com um único 1.

\subsection{Resultado}

A Tabela~\ref{tab:q1-onehot} mostra a matriz. A coluna \texttt{Produto} fica ao lado dos indicadores. Os cabeçalhos omitem o prefixo \texttt{Produto\_}. O arquivo \texttt{scripts/results/q1.csv} guarda as quatro colunas binárias.

\begin{table}[H]
    \centering
    \caption{Codificação one-hot das categorias de produto.}
    \label{tab:q1-onehot}
    \footnotesize
    \setlength{\tabcolsep}{4pt}
    \begin{tabular}{@{}lcccc@{}}
        \toprule
        Produto &
        \shortstack{Eletro-\\domésticos} &
        \shortstack{Eletrô-\\nicos} &
        Livros &
        Roupas \\
        \midrule
        Eletrônicos      & 0 & 1 & 0 & 0 \\
        Roupas           & 0 & 0 & 0 & 1 \\
        Eletrodomésticos & 1 & 0 & 0 & 0 \\
        Eletrônicos      & 0 & 1 & 0 & 0 \\
        Livros           & 0 & 0 & 1 & 0 \\
        Roupas           & 0 & 0 & 0 & 1 \\
        Livros           & 0 & 0 & 1 & 0 \\
        \bottomrule
    \end{tabular}
\end{table}

\section{Questão 2}

\subsection{Enunciado}

Considere o seguinte conjunto de dados:
\begin{verbatim}
data = {
    'Tamanho': [150, 120, 180, 130, 160],
    'Quartos': [3, 2, 4, 3, 4],
    'Preço': [200000, 180000, 250000, 190000, 220000],
    'Localização': [
        'Norte', 'Leste', 'Oeste', 'Sul', 'Centro'
    ]
}
\end{verbatim}
Aplique a técnica de seleção de features para selecionar as features mais relevantes para o negócio.

\subsection{Solução}

O alvo é \texttt{Preço}. As entradas são \texttt{Tamanho}, \texttt{Quartos} e \texttt{Localização}. \texttt{Localização} entra como one-hot. A seleção usa \texttt{SelectKBest} com \texttt{f\_regression} e \(k = 2\): a estatística \(F\) da regressão linear de cada feature com o preço. As duas de maior \(F\) são mantidas.

\subsection{Resultado}

Na amostra, \texttt{Tamanho} (\(F = 54{,}56\)) e \texttt{Quartos} (\(F = 10{,}62\)) apresentaram os maiores F-scores e foram selecionadas pelo \texttt{SelectKBest}. \texttt{Localização} foi codificada em indicadores one-hot e avaliada, mas não entrou entre as duas primeiras. A Tabela~\ref{tab:q2-rank} traz o ranking completo.

\begin{table}[H]
    \centering
    \caption{Ranking das features pela estatística \(F\) em relação ao preço.}
    \label{tab:q2-rank}
    \footnotesize
    \setlength{\tabcolsep}{6pt}
    \begin{tabular}{@{}lccc@{}}
        \toprule
        Feature & \(F\) & \(p\) & Selecionada \\
        \midrule
        Tamanho            & 54,56 & 0,005 & Sim \\
        Quartos            & 10,62 & 0,047 & Sim \\
        Localização\_Oeste &  7,56 & 0,071 & Não \\
        Localização\_Leste &  1,40 & 0,322 & Não \\
        Localização\_Sul   &  0,45 & 0,549 & Não \\
        Localização\_Centro&  0,19 & 0,695 & Não \\
        Localização\_Norte &  0,08 & 0,796 & Não \\
        \bottomrule
    \end{tabular}
\end{table}

\section{Questão 3}

\subsection{Enunciado}

Considere o seguinte conjunto de dados:
\begin{verbatim}
data = {
    'Matemática': [85, 92, 78, 96, 88],
    'Português': [7, 5, 8, 9, 6]
}
\end{verbatim}
É conhecido que as notas de Matemática variam de 0 a 100, enquanto as notas de Português variam de 0 a 10. Aplique a normalização Min-Max e a padronização Z-Score para as notas em ambas as disciplinas, de modo que elas estejam na mesma escala e distribuição.

\subsection{Solução}

As notas de Matemática pertencem à faixa de 0 a 100 e as de Português à faixa de 0 a 10. Os valores brutos, por isso, não compartilham escala. A normalização Min-Max da aula reescala cada coluna para o intervalo \([0, 1]\):
\[
z_{\mathrm{norm}} = \frac{x - x_{\min}}{x_{\max} - x_{\min}}.
\]
\(x_{\min}\) e \(x_{\max}\) são o mínimo e o máximo observados na própria coluna, como no exemplo numérico da aula. A padronização Z-score coloca cada coluna com média zero e desvio padrão unitário:
\[
z = \frac{x - \mu}{\sigma}.
\]
O script ajusta \texttt{MinMaxScaler} e \texttt{StandardScaler}. O desvio é o populacional, o mesmo do exemplo da aula. Em Matemática, o mínimo é 78, o máximo é 96, a média é 87,8 e o desvio é 6,145. Em Português, o mínimo é 5, o máximo é 9, a média é 7 e o desvio é 1,414.

\subsection{Resultado}

Depois das duas transformações, as disciplinas passam a ocupar o intervalo \([0, 1]\) e, na padronização, a mesma distribuição (média 0 e desvio 1). A Tabela~\ref{tab:q3-escalas} reúne os valores. O arquivo \path{scripts/results/q3_escalas.csv} guarda a tabela.

\begin{table}[H]
    \centering
    \caption{Notas originais, normalização Min-Max e padronização Z-score.}
    \label{tab:q3-escalas}
    \footnotesize
    \setlength{\tabcolsep}{4pt}
    \begin{tabular}{@{}rrrrrr@{}}
        \toprule
        \shortstack{Mat.\\orig.} &
        \shortstack{Port.\\orig.} &
        \shortstack{Mat.\\Min-Max} &
        \shortstack{Port.\\Min-Max} &
        \shortstack{Mat.\\Z} &
        \shortstack{Port.\\Z} \\
        \midrule
        85 & 7 & 0,389 & 0,500 & $-0{,}456$ & $0{,}000$ \\
        92 & 5 & 0,778 & 0,000 & $0{,}683$ & $-1{,}414$ \\
        78 & 8 & 0,000 & 0,750 & $-1{,}595$ & $0{,}707$ \\
        96 & 9 & 1,000 & 1,000 & $1{,}334$ & $1{,}414$ \\
        88 & 6 & 0,556 & 0,250 & $0{,}033$ & $-0{,}707$ \\
        \bottomrule
    \end{tabular}
\end{table}

\section{Questão 4}

\subsection{Enunciado}

Considere o seguinte conjunto de dados:
\begin{verbatim}
data = {
    'Nota_Matemática': [8, 7, 9, 6, 8],
    'Nota_Português': [7, 6, 8, 5, 7],
    'Satisfação_Curso': [4, 3, 5, 2, 4]
}
\end{verbatim}
Aplique uma técnica de transformação não linear para criar novas features que capturem as relações entre as notas e a satisfação do curso.

\subsection{Solução}

A aula descreve a transformação não linear como uma operação cujo resultado não é proporcional ao valor original e que torna visível uma relação que as colunas brutas não mostram. Polynomial Features de grau 2 cumpre esse papel: para cada observação, o script gera o quadrado de cada variável e o produto de cada par. Os produtos da nota de Matemática e da nota de Português com a satisfação do curso ligam o desempenho à satisfação. Os quadrados introduzem a curvatura de cada variável. O termo constante fica de fora (\texttt{include\_bias=False}).

\subsection{Resultado}

A Tabela~\ref{tab:q4-poly} mostra as colunas originais e as seis features novas. NM é a nota de Matemática, NP a de Português e Sat. a satisfação do curso. As duas primeiras linhas e a última repetem o mesmo perfil de notas e, por isso, as mesmas features. O arquivo \path{scripts/results/q4_polynomial.csv} guarda a matriz.

\begin{table}[H]
    \centering
    \caption{Features polinomiais de grau 2 para notas e satisfação.}
    \label{tab:q4-poly}
    \footnotesize
    \setlength{\tabcolsep}{2pt}
    \begin{tabular}{@{}ccccccccc@{}}
        \toprule
        NM & NP & Sat. &
        NM\(^2\) &
        \shortstack{NM\(\times\)\\NP} &
        \shortstack{NM\(\times\)\\Sat.} &
        NP\(^2\) &
        \shortstack{NP\(\times\)\\Sat.} &
        Sat.\(^2\) \\
        \midrule
        8 & 7 & 4 & 64 & 56 & 32 & 49 & 28 & 16 \\
        7 & 6 & 3 & 49 & 42 & 21 & 36 & 18 &  9 \\
        9 & 8 & 5 & 81 & 72 & 45 & 64 & 40 & 25 \\
        6 & 5 & 2 & 36 & 30 & 12 & 25 & 10 &  4 \\
        8 & 7 & 4 & 64 & 56 & 32 & 49 & 28 & 16 \\
        \bottomrule
    \end{tabular}
\end{table}

\section{Questão 5}

\subsection{Enunciado}

Considere o seguinte conjunto de dados:
\begin{verbatim}
data = {
    'Idade': [28, 35, 45, 22, 31],
    'Produtos_Comprados': [5, 9, 7, 3, 6]
}
\end{verbatim}
Aplique uma técnica de transformação linear para criar novas features que correlacionem a Idade com Produtos\_Comprados.

\subsection{Solução}

A aula lista a criação de features de interação entre as transformações lineares. A operação é o produto de duas colunas, no mesmo formato de \texttt{valor\_total = preço \(\times\) quantidade}, e serve para correlacionar dois atributos numa coluna só. A feature nova é \texttt{Idade\_x\_Produtos} = Idade \(\times\) Produtos comprados. A correlação de Pearson entre as colunas originais resume a associação linear que esse produto carrega.

\subsection{Resultado}

A correlação de Pearson entre idade e produtos comprados é 0,729: na amostra, idade maior acompanha mais compras. A Tabela~\ref{tab:q5-interacao} traz o produto em cada linha. O arquivo \path{scripts/results/q5_interacao.csv} guarda a tabela.

\begin{table}[H]
    \centering
    \caption{Feature de interação entre idade e produtos comprados.}
    \label{tab:q5-interacao}
    \footnotesize
    \setlength{\tabcolsep}{8pt}
    \begin{tabular}{@{}rrr@{}}
        \toprule
        Idade & Produtos comprados & Idade \(\times\) produtos \\
        \midrule
        28 & 5 & 140 \\
        35 & 9 & 315 \\
        45 & 7 & 315 \\
        22 & 3 &  66 \\
        31 & 6 & 186 \\
        \bottomrule
    \end{tabular}
\end{table}

\section{Questão 6}

\subsection{Enunciado}

Considere o conjunto de dados \texttt{transacoes\_bancarias.csv}, disponibilizado no Canvas:
\begin{itemize}
    \item Aplique uma técnica de Feature Engineering para descobrir correlações entre Sexo e tipo de produto comprado.
    \item Aplique uma técnica de Feature Engineering para identificar qual dia da semana (domingo a sábado) ocorre mais compras.
    \item Aplique uma técnica de Feature Engineering para dimensionar os valores dos produtos por estado numa escala de 0 a 1.
\end{itemize}

\subsection{Sexo e tipo de produto}

\texttt{Sexo} e \texttt{Produto} são categóricos. O script codifica \texttt{Produto} com \texttt{get\_dummies}, como na Questão 1, e \texttt{Sexo} como indicador (Feminino = 1, Masculino = 0). A correlação de Pearson entre esse indicador e cada produto mede a associação. O teste qui-quadrado da tabela de contingência, um dos \textit{filter methods} da aula, avalia se as duas categorias são independentes. O V de Cramér, calculado a partir desse qui-quadrado, resume a intensidade entre 0 e 1.

Nas 1.000 transações, o qui-quadrado é 12,13, com 19 graus de liberdade e \(p = 0{,}880\). O V de Cramér é 0,110. A associação é fraca. A maior correlação absoluta é 0,045 (cafeteira, mais frequente entre mulheres) e todos os \(p\) por produto ficam acima de 0,16. A Tabela~\ref{tab:q6-sexo} lista os vinte produtos. O arquivo \path{scripts/results/q6_sexo_produto.csv} guarda a tabela.

\begin{table}[H]
    \centering
    \caption{Associação entre sexo e tipo de produto. Correlação com o indicador Feminino = 1.}
    \label{tab:q6-sexo}
    \scriptsize
    \setlength{\tabcolsep}{3pt}
    \begin{tabular}{@{}lrrrr@{}}
        \toprule
        Produto & Feminino & Masculino & \(r\) & \(p\) \\
        \midrule
        Cafeteira              & 31 & 21 & $0{,}045$  & 0,166 \\
        Cadeira de escritório  & 32 & 22 & $0{,}044$  & 0,174 \\
        Smartphone             & 21 & 30 & $-0{,}041$ & 0,208 \\
        Notebook               & 27 & 20 & $0{,}033$  & 0,307 \\
        Máquina de lavar       & 28 & 21 & $0{,}032$  & 0,317 \\
        Fogão                  & 17 & 23 & $-0{,}031$ & 0,343 \\
        Console                & 20 & 26 & $-0{,}029$ & 0,376 \\
        Aspirador de pó        & 20 & 25 & $-0{,}024$ & 0,456 \\
        Micro-ondas            & 22 & 27 & $-0{,}023$ & 0,475 \\
        Liquidificador         & 20 & 24 & $-0{,}020$ & 0,546 \\
        Tablet                 & 26 & 22 & $0{,}019$  & 0,564 \\
        Geladeira              & 28 & 24 & $0{,}018$  & 0,579 \\
        Laptop                 & 27 & 31 & $-0{,}017$ & 0,599 \\
        Sofá                   & 26 & 23 & $0{,}014$  & 0,668 \\
        Ventilador             & 28 & 25 & $0{,}013$  & 0,680 \\
        Cama de casal          & 17 & 19 & $-0{,}011$ & 0,739 \\
        Televisão              & 25 & 27 & $-0{,}009$ & 0,782 \\
        Secador de cabelo      & 26 & 28 & $-0{,}009$ & 0,785 \\
        Computador             & 30 & 32 & $-0{,}008$ & 0,799 \\
        Máquina de café        & 29 & 30 & $-0{,}004$ & 0,896 \\
        \bottomrule
    \end{tabular}
\end{table}

\balance
\subsection{Dia da semana}

A aula extrai o dia da semana de uma coluna de data e agrega contagens por grupo. \texttt{Data} é convertida para datetime e cada compra recebe o dia, de domingo a sábado. A contagem por dia indica onde as compras se concentram.

Sábado reúne o maior volume: 163 compras, 16,3\% do total. Segunda-feira reúne o menor, 110. A Tabela~\ref{tab:q6-dias} traz a semana inteira. O arquivo \path{scripts/results/q6_dia_semana.csv} guarda a contagem.

\begin{table}[H]
    \centering
    \caption{Compras por dia da semana.}
    \label{tab:q6-dias}
    \footnotesize
    \setlength{\tabcolsep}{10pt}
    \begin{tabular}{@{}lrr@{}}
        \toprule
        Dia & Compras & Proporção \\
        \midrule
        Domingo       & 131 & 13,1\% \\
        Segunda-feira & 110 & 11,0\% \\
        Terça-feira   & 147 & 14,7\% \\
        Quarta-feira  & 157 & 15,7\% \\
        Quinta-feira  & 144 & 14,4\% \\
        Sexta-feira   & 148 & 14,8\% \\
        Sábado        & 163 & 16,3\% \\
        \bottomrule
    \end{tabular}
\end{table}

\subsection{Valor por estado}

A escala de 0 a 1 é a normalização Min-Max da aula, aplicada dentro de cada estado. O campo \texttt{Estado} chega como o par (UF, nome) e é separado antes do agrupamento. Em cada UF, o menor \texttt{Valor} vai a 0 e o maior vai a 1:
\[
z_{\mathrm{UF}} = \frac{x - x_{\min,\mathrm{UF}}}{x_{\max,\mathrm{UF}} - x_{\min,\mathrm{UF}}}.
\]

Os 27 estados fecham o intervalo: o normalizado mínimo é 0 e o máximo é 1 em todos eles. A Tabela~\ref{tab:q6-amostra} mostra, para três estados, o extremo inferior, um valor intermediário e o extremo superior já na escala nova. A Tabela~\ref{tab:q6-estados} lista a faixa original de cada estado. Os arquivos \path{scripts/results/q6_valor_estado.csv} e \path{scripts/results/q6_valor_norm.csv} guardam o resumo e as 1.000 linhas transformadas.

\begin{table}[H]
    \centering
    \caption{Amostra da escala 0 a 1 do valor, por estado.}
    \label{tab:q6-amostra}
    \footnotesize
    \setlength{\tabcolsep}{4pt}
    \begin{tabular}{@{}llrr@{}}
        \toprule
        UF & Produto & Valor & Escala \\
        \midrule
        AC & Liquidificador &  418,24 & 0,000 \\
        AC & Smartphone     & 5140,19 & 0,495 \\
        AC & Sofá           & 9961,79 & 1,000 \\
        MG & Sofá           & 1178,04 & 0,000 \\
        MG & Notebook       & 5310,23 & 0,490 \\
        MG & Console        & 9617,76 & 1,000 \\
        SP & Computador     &  957,21 & 0,000 \\
        SP & Laptop         & 6017,55 & 0,561 \\
        SP & Ventilador     & 9980,54 & 1,000 \\
        \bottomrule
    \end{tabular}
\end{table}

\begin{table}[H]
    \centering
    \caption{Faixa de valor por estado antes da escala 0 a 1. Em todos, o normalizado vai de 0 a 1.}
    \label{tab:q6-estados}
    \scriptsize
    \setlength{\tabcolsep}{3pt}
    \begin{tabular}{@{}llrrr@{}}
        \toprule
        UF & Estado & \(n\) & Mínimo & Máximo \\
        \midrule
        AC & Acre                  & 40 &  418,24 & 9961,79 \\
        AL & Alagoas               & 27 &  619,77 & 9231,89 \\
        AM & Amazonas              & 42 &  286,91 & 9774,14 \\
        AP & Amapá                 & 34 &  127,90 & 9220,49 \\
        BA & Bahia                 & 37 &   57,47 & 9792,74 \\
        CE & Ceará                 & 36 &  658,82 & 9593,56 \\
        DF & Distrito Federal      & 44 &   74,63 & 9975,50 \\
        ES & Espírito Santo        & 32 &   58,23 & 9562,80 \\
        GO & Goiás                 & 39 &  339,72 & 9546,43 \\
        MA & Maranhão              & 40 &  427,39 & 9857,82 \\
        MG & Minas Gerais          & 39 & 1178,04 & 9617,76 \\
        MS & Mato Grosso do Sul    & 40 &  367,38 & 9961,28 \\
        MT & Mato Grosso           & 42 &   25,42 & 9872,53 \\
        PA & Pará                  & 30 &  194,54 & 9970,01 \\
        PB & Paraíba               & 39 &  193,99 & 9896,82 \\
        PE & Pernambuco            & 23 &  328,53 & 9919,73 \\
        PI & Piauí                 & 32 &  259,19 & 9680,84 \\
        PR & Paraná                & 36 &  474,38 & 9992,45 \\
        RJ & Rio de Janeiro        & 39 &  100,88 & 9874,21 \\
        RN & Rio Grande do Norte   & 36 &  169,31 & 9753,79 \\
        RO & Rondônia              & 49 &  258,74 & 9513,47 \\
        RR & Roraima               & 33 &  301,81 & 9778,12 \\
        RS & Rio Grande do Sul     & 39 &  261,08 & 9957,18 \\
        SC & Santa Catarina        & 29 &  620,77 & 9875,35 \\
        SE & Sergipe               & 41 &   71,92 & 9858,92 \\
        SP & São Paulo             & 39 &  957,21 & 9980,54 \\
        TO & Tocantins             & 43 &  474,47 & 9849,86 \\
        \bottomrule
    \end{tabular}
\end{table}


\end{document}
\endinput
%%
%% End of file `sample-sigconf.tex'.
